We now make precise the sense in which our solver ``knows what it is missing.'' The development is
elementary; we present it as a \emph{formalization} (not a deep theorem) whose value is that it states
exactly---and honestly---what the residual does and does not tell us.

\paragraph{Setup.} A task is $\tau=\{(x_i,y_i)\}_{i=1}^n$ of input--output grids. A \emph{library} $L$ is a
finite set of total functions on grids; a \emph{program} $p=g_{k}\circ\cdots\circ g_{1}$ is a composition of
primitives $g_j\in L$. We assume a \emph{uniform} code over the current library, so a program of length $|p|$
costs $|p|\log_2|L|$ bits (this is a modeling choice; a non-uniform prior changes the constants below). We
score a program by a two-part code:
\begin{equation}
\mathrm{DL}_L(\tau\mid p)\;=\;\underbrace{|p|\log_2|L|}_{\text{program code}}\;+\;\sum_{i=1}^{n}\underbrace{\ell\big(y_i\mid p(x_i)\big)}_{\text{residual code}}.
\label{eq:dl}
\end{equation}
When $\hat y$ and $y$ have the same shape ($N$ cells, $c$ colors) and differ in $d$ cells, the residual code
sends the mismatch count, the mask, and the corrected values, costing
$\log_2(N{+}1)+\log_2\binom{N}{d}+d\log_2 c$ bits (the $\log_2(N{+}1)$ header encodes $d$ and makes the code a
valid prefix code); when shapes differ it encodes $y$ from scratch. Since the header is paid by \emph{every}
hypothesis, it cancels in all comparisons; we therefore define
$\ell(y\mid\hat y)=\log_2\binom{N}{d}+d\log_2 c$ as the \emph{excess} residual cost, which satisfies
$\ell(y\mid\hat y)=0\iff\hat y=y$. Fix any minimizer $p^{\*}_L\in\arg\min_p\mathrm{DL}_L(\tau\mid p)$ and let
the \emph{residual} be $R_L(\tau)=\sum_i \ell(y_i\mid p^{\*}_L(x_i))$. Note the per-symbol cost changes with the
library: adding a primitive makes it $\log_2|L\cup\{g\}|$, and our description-length differences below account
for both program-code terms.

\begin{definition}[Sufficiency]
$L$ is \emph{sufficient} for $\tau$ if some program over $L$ achieves $\sum_i\ell(y_i\mid p(x_i))=0$;
otherwise $L$ is \emph{insufficient}. (Hence ``insufficient'' $\iff R_L(\tau)>0$ by definition.)
\end{definition}

\begin{proposition}[Reuse pays off for frequent-enough reuse]\label{prop:reuse}
Let a program reuse a subroutine $s$ exactly $k$ times. Inlining costs $k|s|\log_2|L|$. Defining $s$ once as a
macro costs (i) the macro definition $|s|\log_2|L|$ plus a table-registration overhead $\beta$, (ii) $k$
references at $\log_2 m$ each, and (iii) an inflation of \emph{every} remaining program symbol from $\log_2|L|$
to $\log_2 m$, where $m=|L|+(\text{\#macros})$. The macro code is shorter once $k$ exceeds a threshold
$k_0$ that grows with $\beta$ and the inflation term; in particular $k=2$ is \emph{not} in general sufficient.
Thus a hierarchical MDL prior prefers compositional explanations \emph{when reuse is frequent enough}---the
expressivity a flat predictor cannot represent, and a partial account of the cliff of \S\ref{sec:cliff}.
\end{proposition}

\begin{theorem}[The residual bounds, localizes, and closes the gap]\label{thm:residual}
Let $\tau$ be generated by a true program $p_{\mathrm{true}}$ over a library $L^{+}$, and let
$L\subsetneq L^{+}$ be insufficient for $\tau$; fix a minimizer $p^{\*}_L$.
\begin{enumerate}[label=(\alph*),leftmargin=1.8em,itemsep=1pt]
  \item \textbf{(Positivity, immediate from Def.~1)} $R_L(\tau)>0$.
  \item \textbf{(Localization)} if $p^{\*}_L$ outputs grids of the correct shape, the residual code charges
  exactly the cells where $p^{\*}_L$ disagrees with $p_{\mathrm{true}}$; otherwise it charges the whole grid.
  Thus the residual localizes \emph{where} a missing operation must act.
  \item \textbf{(Constructive sufficiency / gap-closability)} let $G\subseteq L^{+}\setminus L$ be the set of
  missing primitives that $p_{\mathrm{true}}$ uses ($G\neq\emptyset$, finite). Then $L\cup G$ is sufficient,
  so the minimal residual over $L\cup G$ is $0$ \emph{(unconditional)}. In the \emph{singleton-gap regime}
  $|G|=1$---the case our ablation and cellular-automaton experiments instantiate---a single named primitive
  closes the gap. Moreover the total description length strictly drops,
  $\min_p \mathrm{DL}_{L\cup G}(\tau\mid p)<\mathrm{DL}_L(\tau\mid p^{\*}_L)$, \emph{provided}
  \begin{equation}
  R_L(\tau)\;>\;|p_{\mathrm{true}}|\log_2|L\cup G|\;-\;|p^{\*}_L|\log_2|L|.
  \label{eq:cond}
  \end{equation}
\end{enumerate}
\end{theorem}

\begin{proof}[Proof sketch]
(a) $R_L=0$ would make $L$ sufficient, contradicting the hypothesis. (b) By construction $\ell$ charges
nonzero bits exactly on cells with $p^{\*}_L(x_i)\neq y_i$; since $p_{\mathrm{true}}(x_i)=y_i$ these are the
divergence cells (when shapes match; otherwise the whole grid is charged). (c) Every primitive of
$p_{\mathrm{true}}$ lies in $L\cup G$ by the definition of $G$, so $p_{\mathrm{true}}$ is expressible over
$L\cup G$ and the minimal residual there is $0$; when $|G|=1$ this is the addition of a single primitive.
(Note a single $g$ does \emph{not} suffice in general: if $p_{\mathrm{true}}$ composes two missing
primitives, no singleton extension is sufficient---hence the explicit set $G$.) For the strict DL drop, take
$p'=p_{\mathrm{true}}$: then $\min_p\mathrm{DL}_{L\cup G}(\tau\mid p)\le|p_{\mathrm{true}}|\log_2|L\cup G|$,
while $\mathrm{DL}_L(\tau\mid p^{\*}_L)=|p^{\*}_L|\log_2|L|+R_L$; subtracting gives the drop, which is positive
exactly under~\eqref{eq:cond}.
\end{proof}

\paragraph{What this does and does not say.} Part (c) is named \emph{gap-closability}, not identifiability,
deliberately: the theorem guarantees that re-adding the missing primitive set $G$ (a \emph{single} primitive
in the singleton-gap regime our experiments instantiate) drives the residual to zero and
(under~\eqref{eq:cond}) lowers description length---but it does \emph{not} prove that the residual
\emph{uniquely identifies} the missing primitive. Several primitives may close a gap, and partial compensation by other
primitives can blur the residual's structural pattern. We therefore distinguish two diagnostic tiers and
measure both: a cheap \emph{residual-signature classifier} (``what category of operation is missing?'') and a
\emph{residual-guided re-add search} (``which primitive closes it?''). On a controlled domain with known ground
truth (\S\ref{sec:generality}), the re-add tier closes $100\%$ of provably-needy gaps---the constructive
content of (c)---while the cheap signature recovers the correct category only $62\%$ of the time. We thus do
\emph{not} claim the residual is ``exactly the signature'' of the missing primitive; we claim it localizes the
gap, certifies closability in the singleton regime, and serves as an empirically useful but imperfect cue for
\emph{naming} what is missing.
